\subsection{All marked simple even bases and their relative pronilpotent kernels}

\begin{statement}{D.43}{Proved}{A simultaneous full-family extension and an evaluated orbit of two simple factors}
Let $B_{\rm even}=\ker(B\to C_2)$, with the original
sign map, and consider the entire marked family
\[
 \mathscr J=\{N\lhd B:\ N\subset B_{\rm even},\
       B_{\rm even}/N\text{ is finite nonabelian simple}\}.
                                                               \tag{D.43.1}
\]
There is one coordinate for each actual source
kernel. Write $L_N=B_{\rm even}/N$ and $J_N=B/N$.
The actual image $s_N$ gives
$J_N=L_N\rtimes\langle s_N\rangle$.
For each prime $p$, put
\[
 R_{N,p}=\ker(\widehat N\to N^{(p)}),\qquad
 \mathcal P_N=\widehat B/\bigcap_pR_{N,p},\qquad
 R_{\mathscr J}=\bigcap_{N\in\mathscr J,\ p}R_{N,p}.
                                                               \tag{D.43.2}
\]
All completions and intersections are taken with
these indicated source maps. Each $N$ is free
of rank $1+|L_N|/3$, as its marked
Bass--Serre graph has $2|L_N|$ edges and
$|L_N|+2|L_N|/3$ vertices. Their actual
simultaneous images are
\[
 \begin{split}
 J_{\mathscr J}
    &=\left(\prod_{N\in\mathscr J}L_N\right)\rtimes C_2,\\
 \mathcal P_{\mathscr J}
    :=\widehat B/R_{\mathscr J}
    &=\mathop{\prod_{\mathcal R_3}}_{N\in\mathscr J}
                                      \mathcal P_N ,
 \end{split}                                                   \tag{D.43.3}
\]
where $\mathcal R_3$ is the exact relative pro-$3$
completion (D.42.13). In the first line the
involution acts on each coordinate by its actual
conjugation action; it does not permute the
coordinates. The maximal prosolvable quotient
of $\mathcal P_{\mathscr J}$ is exactly
$\mathcal R_3$.

Here is an exact common-quotient formula that
retains the perfect and symmetric coordinates
already present in D.41. Put
\[
 \begin{split}
 \mathcal D&=\mathcal R\times_{C_2}\mathcal S
                    =\widehat B/K_D,\qquad
 K_D=K_{\mathcal R}\cap K_{\mathcal S},\\
 \mathcal D_{\mathscr J}
    &=\widehat B/(K_D R_{\mathscr J}),\qquad
 \mathcal C_{\mathscr J}
    =\mathcal C_{54}\times_{C_6}\mathcal D_{\mathscr J}.
 \end{split}                                                   \tag{D.43.4}
\]
The products of closed normal subgroups in this
display are closed. The maps to $C_6$ are the
original full abelianization of $B$.
Then $\mathcal C_{\mathscr J}$ is the maximal
common marked quotient of $\mathcal W$ and
$\mathcal P_{\mathscr J}$, with mutually inverse
maps proved below. In particular
\[
 \mathcal Z_{43}
   :=\widehat B/(R_W\cap R_{\mathscr J})
    =\mathcal W\times_{\mathcal C_{\mathscr J}}
                                \mathcal P_{\mathscr J}.        \tag{D.43.5}
\]
The automorphism
\[
       (\alpha_{41},1_{\mathcal P_{\mathscr J}})                 \tag{D.43.6}
\]
is an actual continuous peripheral automorphism
of (D.43.5), with multiplier one, the same
global gauge, and both hexagons. One simultaneous
centralizer lift supplies its conjugators.
It retains D.42 exactly as an action under the
canonical projection and acts as the identity
on every relative pro-$p$ completion in
(D.43.2), for all bases and all primes at once.

All three finite-base types enter this formula.
For the inner-action bases, the proof gives
the exact isomorphism $J_N\to L_N\times C_2$
and its actual perfect-coordinate map into
$\mathcal R$. Full symmetric bases use their
actual coordinates in $\mathcal S$. Let
$\mathscr J_{\rm new}$ consist of the remaining
bases: the involution has a non-inner action
on $L_N$, and $J_N$ is not a full symmetric
group. For this entire subfamily the formula
simplifies to
\[
 \mathcal D_{\mathscr J_{\rm new}}=C_6,\qquad
 \mathcal C_{\mathscr J_{\rm new}}=\mathcal C_{54}.
                                                               \tag{D.43.7}
\]
It includes both marked projective-linear
groups in D.41 and D.42. For the whole
family, (D.43.4) specifies its additional
properfect common part by its exact source
kernel; it has not been replaced by a
sign quotient. Its commutator subgroup is
properfect, and it maps onto the simultaneous
image of every retained inner and symmetric
base.

The simultaneous pronilpotent base kernel
also has an explicit form. The source
inclusions give surjections
$N^{(3)}\to B_{\rm even}^{(3)}$. With
$\mathcal N=\ker(\mathcal P_{\mathscr J}
                       \to J_{\mathscr J})$,
\[
 \mathcal N
   =\left(\prod_{\substack{N\in\mathscr J\\p\ne3}}
                                       N^{(p)}\right)
       \times
        \left(\mathop{\prod_{B_{\rm even}^{(3)}}}_{N\in\mathscr J}
                                       N^{(3)}\right).          \tag{D.43.8}
\]
In particular $\mathcal N$ is pronilpotent.
The second factor is its pro-$3$ Sylow
subgroup. If $\delta:\widehat B\to\mathcal D$
is the actual source map, the two comparison
kernels in (D.43.5) are
\[
 \begin{split}
 \ker(\mathcal Z_{43}\to\mathcal P_{\mathscr J})
    &=K_{T,C}\times\delta(R_{\mathscr J}),\\
 \ker(\mathcal Z_{43}\to\mathcal W)
    &=\ker(\mathcal P_{\mathscr J}\to\mathcal C_{54})
        \cap
      \ker(\mathcal P_{\mathscr J}\to\mathcal D_{\mathscr J}).
 \end{split}                                                   \tag{D.43.9}
\]
$K_{T,C}$ and all its successive kernels
are exactly those of (D.42.9). The second
line is an intersection of the displayed
actual quotient maps inside
$\mathcal P_{\mathscr J}$.

A further explicit finite source quotient
remains. On the symbols $1,\ldots,7$, put
\[
 a=(1\ 4\ 7),\qquad
 b=(1\ 2\ 3)(4\ 5\ 6),\qquad
 J_*=(A_7\times A_7)\rtimes\langle\iota\rangle ,
                                                               \tag{D.43.10}
\]
where $\iota$ interchanges the two factors.
There is an actual epimorphism
\[
 \sigma_*:\widehat B\twoheadrightarrow J_*,
       \qquad s\longmapsto\iota,\quad r\longmapsto(a,b).
                                                               \tag{D.43.11}
\]
Its even image has order $6350400$, its
whole image has order $12700800$, and
the actual orders of $t,x$ are $14,7$.
The remaining full-source kernel has
the exact image
\[
       \sigma_*(R_W\cap R_{\mathscr J})
                         =A_7\times A_7.                       \tag{D.43.12}
\]
Thus the new completion still has a nontrivial
kernel in $\widehat B$, and likewise in the
indicated $\widehat F$ source. This calculation
does not exclude the retained CM pair from
a full profinite lift. D.45 excludes the particular recorded
action extended here and leaves the CM pair as a target
for a replacement.

This addresses Part D(e), section 51.16,
folio 734 verso, scan 174, and
D-after-LM2(1), section 51.15,
folios 725--727, scans 156--160.
\end{statement}
\begin{proof}
\emph{The entire marked family and its exact simultaneous images.}
For each $N$, the restriction
$B_{\rm even}\to L_N$ has kernel exactly
$N$. Distinct members of (D.43.1)
therefore give distinct marked simple
quotients of the same even source.
The simple subdirect-product argument
of D.41 gives surjectivity onto every
finite product of them: a failure on
adjoining one factor would identify
its kernel with that of a preceding
factor. Compactness gives the full
product. The actual image of $s$
then proves the first line of
(D.43.3), with its stated action.
Repeated epimorphisms having the same
source kernel differ by the unique
induced isomorphism of their targets.
Their images are its graph; deletion
and reinsertion using that isomorphism
are inverse marked maps. The same
source kernel in (D.43.2) gives the
identical relative completion, so
this also treats repetitions at
every prime and every finite level.

The finite kernel of a simultaneous
relative-pronilpotent image over
a finite joint base embeds in a
product of nilpotent kernels.
It is therefore nilpotent. Its
even base is a product of
nonabelian simple groups and
is perfect. The proof of
(D.42.15) consequently applies
verbatim to the even source:
after its perfect residual is
killed, every solvable image
is an image of that nilpotent
kernel and is generated by
$r,srs$. It is a $3$-group.
The indicated $\mathcal R_3$
is a quotient of every
$\mathcal P_N$, so it is
exactly the maximal
prosolvable quotient of
every finite-family completion.

For precision, the common quotient
when adjoining a new member $N$
is also exactly $\mathcal R_3$.
Suppose a finite common image
were nonsolvable. Viewed as a
quotient of the new relative
completion, its even subgroup,
modulo the image of the
nilpotent kernel, has the
actual quotient $L_N$.
Indeed the only possible
even images in the simple
base are $1$ and $L_N$;
the first would make the
whole image solvable.
Let $N_I=\bigcap_{M\in I}M$
for the preceding finite
family. The simultaneous
simple-base surjectivity
just proved gives
$N_I\twoheadrightarrow L_N$.
But in the preceding
relative completion the
image of $N_I$ lies in its
pronilpotent base kernel.
Its image in the finite
simple group would be
nilpotent, a contradiction.
Thus every common image
is solvable and the
maximal common quotient
is the stated $\mathcal R_3$.
Induction and compactness
prove the second line of
(D.43.3). This argument
includes inner and symmetric
bases; it never assumes
that the involution's
action is outer.

The actual joint image of
$J_N$ and $\mathcal R_3$
is $J_N\times_{C_2}\mathcal R_3$:
the maximal solvable quotient
of $J_N$ is $C_2$.
Consequently its kernel over
$J_N$ surjects onto
$B_{\rm even}^{(3)}$.
At primes other than three
the map to this pro-$3$
group is trivial. Reordering
the literal fibre-product
coordinates proves (D.43.8).
It also proves that
$\mathcal N$ projects onto
the base kernel of every
finite-family completion:
extend its common
$B_{\rm even}^{(3)}$
coordinate to each omitted
factor using the just-proved
surjections. This projection
property will be used in
the common-quotient and
centralizer arguments.

\emph{The three base types and their actual existing maps.}
Let $\theta_N$ be conjugation
by $s_N$ on $L_N$. When
$\theta_N=\operatorname{Int}(h)$,
the element $h$ is unique
because $L_N$ is centreless,
and $h^2=1$ because
$\theta_N^2=1$. The maps
\[
 \begin{aligned}
 L_N\rtimes\langle s_N\rangle
       &\longrightarrow L_N\times C_2,\\
       l s_N^\epsilon&\longmapsto(lh^\epsilon,\epsilon),\\
 (m,\epsilon)&\longmapsto m h^{-\epsilon}s_N^\epsilon
 \end{aligned}                                                 \tag{D.43.13}
\]
are inverse homomorphisms.
The projection to $L_N$
is the actual perfect
source map
$s\mapsto h,\ r\mapsto r_N$.
It is onto, since
$B\to J_N$ is onto.
Thus this base occurs in
$\mathcal R$ together
with the common sign.
When $J_N$ is a full
$S_n$, its actual kernel
is one of the coordinates
of $\mathscr S$, so it
occurs in $\mathcal S$.
These are the marked
maps used in (D.43.4).

For the remaining bases,
$\theta_N$ is non-inner
and $J_N$ is not a full
symmetric group. Their
only normal subgroups
are $1,L_N,J_N$. To see
this, a normal subgroup
contained in $L_N$ is
trivial or all of $L_N$.
An odd normal subgroup
with trivial intersection
with $L_N$ would centralize
$L_N$ and make
$\theta_N$ inner, which
was excluded. The
finite coordinate-kernel
argument of D.41 now
excludes each such
$J_N$ from $\mathcal D$
and from $\mathcal W$.
Perfect-coordinate
kernels have odd image
and hence image $J_N$;
symmetric-coordinate
kernels have image
containing $L_N$, since
$J_N$ is not a quotient
of a full $S_n$.
The solvable coordinate,
when present, has kernel
image containing $L_N$.
Their iterated commutator
is trivial in the
subdirect source image
but contains the
nontrivial perfect
$L_N$ in the target.

For a finite joint base
in this new subfamily,
the even subgroup is
a product of its simple
factors and the involution
acts non-innerly on each.
Every proper normal
subgroup is a product
of even factors. An
odd normal element
commuted with each
factor would otherwise
force all those factors,
and then the whole group,
into the normal subgroup.
It follows that every
nonsolvable quotient
of a relative
pronilpotent image
retains one of these
actual $J_N$ quotients.
Thus its common quotient
with $\mathcal D$ is
solvable. The maximal
solvable quotient of
$\mathcal D$ is $C_6$,
and both sides retain
that marked quotient.
This proves (D.43.7).
For the whole family,
the maps (D.43.13)
and the symmetric
coordinates instead
show why the properfect
part in (D.43.4)
must be retained.

\emph{The key common-quotient calculation for the whole family.}
Let $Q$ be any finite common
image of $\mathcal W$ and
$\mathcal P_{\mathscr J}$.
Write
\[
 P=\operatorname{im}\left(
      \mathcal R'\times\prod A_M\longrightarrow Q\right),
 \qquad
 Z=\operatorname{im}\left(
      \mathcal T'\longrightarrow\mathcal W
                               \longrightarrow Q\right).
                                                               \tag{D.43.14}
\]
The second map is the
\emph{isolated-coordinate}
map $t\mapsto(t,1,1)$.
It is defined on
$\mathcal T'=\ker(\mathcal T\to C_6)$.
The group $P$ is perfect,
and $Q/P$ is solvable.
Hence $P=Q^\infty$,
the stable derived
subgroup. The subgroup
$Z$ is normal and
centralizes $P$.

Factor the map to $Q$
through a finite-family
relative image retaining
its finite joint base.
Let $V$ be the image
of its nilpotent base
kernel in $Q$.
The even subgroup of
$Q/V$ is a quotient
of a product of
nonabelian simple
groups, hence is
a centreless product
of the surviving
simple factors.
The image of $P$
surjects onto that
even subgroup.
The image of $Z$
is even and centralizes
it, so is trivial.
Thus
\[
                    Z\subset V.                              \tag{D.43.15}
\]
In particular $Z$
is nilpotent.

For a prime $p\ne3$,
let $Z_p$ be its
characteristic Sylow
subgroup. The quotient
$Q/P$ is a relative
$3$-group over the
sign, by the first
part of the proof.
Since $Z_p$ is even,
its image there is
trivial. Thus
$Z_p\subset P$ and
$Z_p$ is central in
$P$. For every
$q\in Q_{\rm even}$,
choose $a\in P$
with the same image
in $(Q/V)_{\rm even}$.
Then $q=av$ for
$v\in V$, and
conjugation by $q$
on $Z_p$ equals
conjugation by $v$.
Because $V$ is
nilpotent, this
action lies in a
$p$-group: all
its other Sylow
subgroups commute
with $Z_p$.
But $Q_{\rm even}$
is generated by
the actual images
of $r,srs$, both
of order dividing
three. Therefore
it acts trivially
on $Z_p$.

Here the final
vanishing argument
uses no projection
from $B_{\rm even}$
to $\mathcal T'$.
Let $B'=\ker(B\to C_6)$.
The actual restriction
$B'\to\mathcal T'$
is onto after
profinite completion,
and its image is
onto the finite
quotient under
consideration.
Compose this map
with (D.43.14)
and the characteristic
projection $Z\to Z_p$:
\[
              \phi:B'\longrightarrow Z_p.                     \tag{D.43.16}
\]
It is invariant
under conjugation
by $B_{\rm even}$,
because $Z_p$
centralizes
$Q_{\rm even}$.
Consequently
$\ker\phi$ is
normal in
$B_{\rm even}$,
and
\[
 1\longrightarrow\operatorname{im}\phi
  \longrightarrow B_{\rm even}/\ker\phi
  \longrightarrow B_{\rm even}/B'=C_3
  \longrightarrow1                                           \tag{D.43.17}
\]
is a central extension.
A group with central
kernel and cyclic
quotient is abelian.
The group in
(D.43.17) is generated
by the actual two
order-three elements
$r,srs$, so it is
abelian of exponent
dividing three.
Its $p$-subgroup
$\operatorname{im}\phi$
is trivial for
$p\ne3$.
The map in
(D.43.16) is onto
$Z_p$, so $Z_p=1$.
This proves that
$Z$ is a $3$-group.

Let $\phi_T:\mathcal T'\to Z$
be the map in
(D.43.14). Its
kernel is normal
in $\mathcal T$:
conjugation can
be calculated using
any lift to
$\mathcal W$, and
two lifts differ
by $P$, which
centralizes $Z$.
The quotient
$\mathcal T/\ker\phi_T$
has normal $3$-group
$Z$ and quotient
$C_6$. Its even
subgroup is therefore
a $3$-group.
It is an actual
relative pro-$3$
quotient over
the sign, so by
D.42 its source
map factors through
$\mathcal C_{54}$.
We have proved
that every finite
common image kills
the isolated subgroup
$K_{T,C}$.

\emph{The exact marked fibre-product formula.}
Let $K_C=\ker(\widehat B\to\mathcal C_{54})$.
The preceding argument
proves that the maximal
common quotient of
$\mathcal W$ and
$\mathcal P_{\mathscr J}$
factors through
\[
 \mathcal W/K_{T,C}
      =\mathcal C_{54}\times_{C_6}\mathcal D
      =\widehat B/(K_C\cap K_D).                       \tag{D.43.18}
\]
Conversely both sources
retain $\mathcal C_{54}$,
so $R_{\mathscr J}\subset K_C$.
Products of closed
normal subgroups are
closed, and the elementary
modular identity gives
\[
 (K_C\cap K_D)R_{\mathscr J}
     =K_C\cap(K_D R_{\mathscr J}).
                                                               \tag{D.43.19}
\]
The common source
kernel is
$R_W R_{\mathscr J}$.
The proved inclusion
from (D.43.18) and
$R_W\subset K_C\cap K_D$
give
\[
 R_W R_{\mathscr J}
   =(K_C\cap K_D)R_{\mathscr J}
   =K_C\cap(K_D R_{\mathscr J}).                        \tag{D.43.20}
\]
These are both
directions of the
comparison, with
the actual source
kernels.

The maximal
prosolvable quotient
of $\mathcal D$
is $C_6$, because
its commutator
subgroup is the
properfect product
$\mathcal R'\times
\prod A_M$.
Both $\mathcal D$
and $\mathcal P_{\mathscr J}$
retain $C_6$,
so their common
quotient
$\mathcal D_{\mathscr J}$
does also, and
its commutator
subgroup is
properfect.
Its maximal common
quotient with
$\mathcal C_{54}$
is therefore
$C_6$.
Closed Goursat
applied to
(D.43.20) proves
(D.43.4)--(D.43.5).
The inner and
symmetric bases
are actual common
quotients of
$\mathcal D$ and
$\mathcal P_{\mathscr J}$;
their simultaneous
image consequently
is a quotient of
$\mathcal D_{\mathscr J}$.
The description
by $K_D R_{\mathscr J}$
keeps any further
properfect common
kernel data as well.
It is not an
assertion that this
group consists
only of its finite
simple bases.

\emph{One simultaneous fixed-element lift.}
On (D.43.18), the
action $\alpha_{41}$
is the identity:
it is the identity
on $\mathcal C_{54}$
by D.42 and on
$\mathcal D$ by
D.40--D.41.
It is consequently
the identity on
$\mathcal C_{\mathscr J}$.
Thus (D.43.6)
preserves the
actual fibre product
and is an
automorphism.
For the peripheral
conjugators one
must also retain
the nonzero
coordinate of
D.42. In
$\mathcal C_{\mathscr J}$
the recorded
conjugators have
images
\[
          A_C=1,\qquad
          \beta_C=(b_C,1),\qquad
          b_C=-\tfrac12(e_7+e_8).                     \tag{D.43.21}
\]
This element
centralizes $s_C$.

Set $V=\operatorname{im}(\mathcal N
\to\mathcal C_{\mathscr J})$.
It is a closed
pronilpotent
normal subgroup.
The projection
property following
(D.43.8) shows
that its image
in every finite
common quotient
is exactly the
image of the
finite-family
base kernel
used in
(D.43.15).
Consequently
the image of
the isolated
$\mathcal T'$
coordinate lies
in $V$.
It is pro-$3$
by the same
finite argument.
In particular
$\beta_C\in V_3$,
the characteristic
pro-$3$ Sylow
subgroup.

The second factor
of (D.43.8)
maps onto $V_3$.
This map is
$s$-equivariant
and has pro-$3$
kernel. Apply
the coprime
involution argument
proved in D.42
to obtain one
element
\[
 \widetilde\beta\in\mathcal N_3,\qquad
 \widetilde\beta\longmapsto\beta_C,\qquad
 s\widetilde\beta s^{-1}=\widetilde\beta.
                                                               \tag{D.43.22}
\]
Explicitly, use
finite quotients
of $\mathcal N_3$
normal in
$\mathcal N_3\rtimes\langle s\rangle$,
together with the
corresponding finite
quotients of $V_3$.
For every finite
set of such
conditions, their
common refinement
has a fixed lift
by the finite
centre-induction
argument. The
closed sets of
solutions in
$\mathcal N_3$
have the finite
intersection
property.
Compactness gives
(D.43.22) at
all bases, primes
and finite levels
simultaneously.

Now the actual
compatible pairs
\[
           (A_W,1),\qquad
           (\beta_W,\widetilde\beta)                  \tag{D.43.23}
\]
belong to
$\mathcal Z_{43}$.
They conjugate
the indicated
images of $r,s$
as required,
and the product
automorphism
fixes $t$.
Its parameter
is
$(f_{39},1)$.
It is the image
of an actual
element of
$\widehat F'$,
since the
$\mathcal W$
coordinate
retains the
whole abelianization
of $\widehat F$
and has zero
coordinate there.
The same
order-two,
order-three
and fixed-cusp
equations give
both multiplier-one
hexagons.
No additional
finite-level
gauge is introduced.
The canonical
projection to
D.42 retains its
entire action;
the possible
choice of a
different centralizing
representative
does not change
that action or
its commutator
parameter.

The formula
(D.43.9) follows
by taking the
two kernels in
the now-proved
fibre product.
For the first,
the map
$\mathcal W=
\mathcal T\times_{C_6}\mathcal D$
to
$\mathcal C_{54}\times_{C_6}
\mathcal D_{\mathscr J}$
has kernel the
direct product
of its two
displayed kernels.
For the second,
the two coordinate
conditions are
exactly those in
(D.43.4). This
proves the
comparison maps
and kernels
without identifying
the additional
properfect part
with a smaller
finite base.

\emph{The evaluated orbit of two simple factors.}
The program enumerates
the subgroup generated
by $a,b$ in (D.43.10)
and obtains all
$2520$ elements
of $A_7$.
Their centralizer
orders are respectively
$36$ and $9$,
so no automorphism
of $A_7$ interchanges
$a$ and $b$.
The two even-source
markings
\[
 \begin{array}{c|cc}
       &r&srs\\ \hline
  \rho_1&a&b\\
  \rho_2&b&a
 \end{array}                                                   \tag{D.43.24}
\]
therefore have
different kernels.
Simple Goursat
gives the full
$A_7^2$ image.
There is also
a direct word
certificate.
The program finds
a $16$-letter
word in
$r,srs$ and
their inverses,
equal to a
$32$-letter word
in $s,r,r^{-1}$,
whose paired
image is
$(1,z)$ with
$z\ne1$.
It enumerates
the $360$
conjugates of
$z$ and verifies
that they generate
all $2520$
elements.
Thus the paired
image contains
$1\times A_7$
and surjects
the first
factor, proving
the same full
image directly.
Conjugation by
$s$ supplies
the other factor.
This verifies
(D.43.11),
all stated orders,
and the faithful
degree-fourteen
permutation
marking.

The normal
subgroups of
$J_*$ are
$1,A_7^2,J_*$.
Indeed a normal
subgroup contained
in $A_7^2$
is a product
of simple
factors, and
the involution
interchanges
them. An odd
normal element
does not
centralize
$A_7^2$;
its commutators
force the whole
even subgroup
and then
$J_*$ into
that normal
subgroup.
The even
subgroup is
perfect.

There is no
$J_*$ quotient
of $\mathcal W$,
by the finite
coordinate-kernel
argument already
used in D.41:
the solvable
coordinate
kernel has
image containing
$A_7^2$,
the perfect
coordinate
kernels have
odd image and
hence image
$J_*$, and
the full
symmetric
coordinate
kernels have
image containing
$A_7^2$.
The latter
assertion uses
the normal
subgroups of
$S_n$: $J_*$
has order
$12700800$,
is not a
full symmetric
group, and
is not a
quotient of
one. The
iterated
commutator
of the
coordinate
kernels
would be
trivial with
nontrivial
perfect image.

Nor is
$J_*$
a quotient
of
$\mathcal P_{\mathscr J}$.
In a finite
relative image
the nilpotent
base kernel
must have
trivial image
in $J_*$:
neither of
its two
nontrivial
normal subgroups
is nilpotent.
The map would
therefore factor
through a finite
joint base in
(D.43.3).
Each of its
simple even
factors is
normal in the
whole base,
because its
involution
preserves
that factor.
Its image
in $J_*$
would be
a simple
normal subgroup.
No such
nontrivial
subgroup
exists in
$J_*$.
All even
factors are
therefore
killed, and
the remaining
$C_2$
cannot surject
onto $J_*$.

Finally a finite
common subdirect
image of
$\mathcal W$
and
$\mathcal P_{\mathscr J}$
cannot map
onto $J_*$
either.
Its two
coordinate
kernels commute;
their images
would both
contain
$A_7^2$,
since neither
coordinate
admits that
quotient.
This contradicts
the nontrivial
perfectness
of $A_7^2$.
Thus
$\mathcal Z_{43}$
has no such
quotient.
Its remaining
source kernel
lies in
$\widehat B'$
and has
normal image
in $J_*$.
The image
cannot be
trivial,
as otherwise
the displayed
epimorphism
would factor
through
$\mathcal Z_{43}$.
The normal
subgroup
classification
gives exactly
(D.43.12).

The independent
standard-library
certificate is
\texttt{d\_cm\_simple\_even\_family.py},
with its saved
output and
JSON word
certificate.
The remaining
kernel in
(D.43.12)
shows explicitly
why no
cofinality
assertion is
being used
for the full
free profinite
source.
\end{proof}
